(1) 对输入向量应用单位分束矩阵，\(U_{\rm BS}(E_0,0)^T=(E_0/\sqrt2,iE_0/\sqrt2)^T\)。输出总功率仍为 \(|E_0|^2\)，因为矩阵是酉矩阵，第二臂振幅带有反射引起的相位因子 \(i\)。
(2) 传播后的双臂向量为 \((E_0e^{i\phi_1}/\sqrt2,iE_0e^{i\phi_2}/\sqrt2)^T\)。再经同一理想分束器，两个输出为 \(E_a=E_0(e^{i\phi_1}-e^{i\phi_2})/2\)、\(E_b=iE_0(e^{i\phi_1}+e^{i\phi_2})/2\)，可直接由矩阵乘法验证。
(3) 令相位差 \(\Delta\phi=\phi_1-\phi_2\)，两端强度分别为 \(I_a/I_0=|e^{i\phi_1}-e^{i\phi_2}|^2/4=\sin^2(\Delta\phi/2)\)、\(I_b/I_0=\cos^2(\Delta\phi/2)\)。因 \(\sin^2+\cos^2=1\)，两端总输出功率恒为输入功率，体现器件无吸收。
(4) 光程差变化 \(\Delta\ell\) 带来 \(\Delta\phi=2\pi\Delta\ell/\lambda_0\)，输出光强随 \(\Delta\ell\) 变化一个真空波长 \(\lambda_0\) 后重复。在完全同相 \(\Delta\ell=0\) 且无其他相移时，\(I_a=0,I_b=I_0\)，体现两个端口的互补干涉。
